{\parskip=0pt\centering\large\textbf{ABSTRAK}\par}

\addcontentsline{toc}{chapter}{ABSTRAK}
\vspace{3ex}

{\parskip=0pt\centering\textbf{\tatitle{}}\par}
\vspace{\baselineskip}

\begingroup
\setlength{\tabcolsep}{0pt}

\bfseries
\begin{tabularx}{\textwidth}{l >{\centering}m{2em} X}
  Nama Mahasiswa / NRP & : & \name{} / \nrp{} \\[1ex]
  Departemen           & : & \department{} \facultyshort{} - ITS \\[1ex]
  Dosen Pembimbing     & : & Nama Pembimbing dan Gelar \\
\end{tabularx}
\endgroup
\vspace{\baselineskip}

\paragraphtextstyle

\noindent\textbf{Abstrak}

% Ubah paragraf berikut dengan abstrak dari proposal tugas akhir
Pada penelitian ini diajukan \lipsum[1]

\vspace{\baselineskip}

\noindent\textbf{Kata kunci:} \textbf{\textit{kata kunci 1, kata kunci 2, kata kunci 3, kata kunci 4, kata kunci 5}}.

